% =====================================================================
%  Preprocesamiento y extracción de características de señales ECG
%  para la detección del síndrome de Brugada
%  Inteligencia de Negocio · Trabajo de teoría 1 · Primera parte del TFG
%  Compilar:  latexmk -pdf main.tex      (pdfLaTeX + biber)
% =====================================================================
\documentclass[11pt, a4paper, oneside, parskip=half, numbers=noenddot,
               listof=totoc, bibliography=totoc, captions=tableheading,
               headings=optiontohead, toc=flat]{scrreprt}
\usepackage{estilo}
\addbibresource{referencias.bib}

\hypersetup{pdftitle={Preprocesamiento y extracción de características de señales ECG para la detección del síndrome de Brugada},
            pdfauthor={César Capelo}, pdfsubject={Inteligencia de Negocio · Trabajo de Fin de Grado},
            pdfkeywords={síndrome de Brugada, ECG, preprocesamiento, selección de características, ETL}}

\begin{document}
\pagenumbering{roman}
% ---------------------------------------------------------------------
%  Portada
% ---------------------------------------------------------------------
\begin{titlepage}
\begin{tikzpicture}[remember picture, overlay]
  % banda superior
  \fill[acento] (current page.north west) rectangle ([yshift=-6.2cm]current page.north east);
  % trazado de ECG estilizado (latido con patrón tipo 1 en el tercer complejo)
  \begin{scope}[shift={([yshift=-4.65cm]current page.north west)}, x=1cm, y=1cm]
    \draw[white, opacity=0.9, line width=1.1pt, line join=round]
      (0,0) -- (1.6,0) .. controls (1.8,0.25) and (2.0,0.25) .. (2.2,0) -- (2.5,0)
      -- (2.62,-0.25) -- (2.8,1.55) -- (2.98,-0.55) -- (3.1,0) -- (3.5,0)
      .. controls (3.8,0.45) and (4.2,0.45) .. (4.5,0) -- (6.0,0)
      .. controls (6.2,0.25) and (6.4,0.25) .. (6.6,0) -- (6.9,0)
      -- (7.02,-0.25) -- (7.2,1.55) -- (7.38,-0.55) -- (7.5,0) -- (7.9,0)
      .. controls (8.2,0.45) and (8.6,0.45) .. (8.9,0) -- (10.4,0)
      .. controls (10.6,0.25) and (10.8,0.25) .. (11.0,0) -- (11.3,0)
      -- (11.42,-0.25) -- (11.55,1.2) -- (11.7,-0.35) -- (11.82,0.75)
      .. controls (12.2,0.62) and (12.6,0.2) .. (12.9,-0.32)
      .. controls (13.1,-0.55) and (13.35,-0.3) .. (13.6,0) -- (15.0,0)
      .. controls (15.2,0.25) and (15.4,0.25) .. (15.6,0) -- (15.9,0)
      -- (16.02,-0.25) -- (16.2,1.55) -- (16.38,-0.55) -- (16.5,0) -- (16.9,0)
      .. controls (17.2,0.45) and (17.6,0.45) .. (17.9,0) -- (21.0,0);
  \end{scope}
  \node[anchor=north west, text=white, font=\sffamily\bfseries\large, inner sep=0]
       at ([xshift=2.5cm, yshift=-1.6cm]current page.north west) {UNIVERSIDAD DE OVIEDO};
  \node[anchor=north west, text=white, font=\sffamily, inner sep=0, opacity=0.9]
       at ([xshift=2.5cm, yshift=-2.25cm]current page.north west) {Escuela de Ingeniería Informática};
\end{tikzpicture}

\vspace*{5.6cm}
{\sffamily\color{acento}
{\small\bfseries\MakeUppercase{Inteligencia de Negocio}\enspace·\enspace Trabajo de teoría 1}\par
\vspace{0.3cm}
{\small\color{gristexto} Primera parte de la memoria del Trabajo de Fin de Grado}\par
\vspace{1.0cm}
{\fontsize{27}{32}\selectfont\bfseries Preprocesamiento y extracción de\\ características de señales ECG\\ para la detección del síndrome\\ de Brugada\par}
\vspace{0.6cm}
{\color{acento!40}\rule{0.32\linewidth}{2pt}}\par
\vspace{0.6cm}
{\large\color{gristexto}\normalfont\sffamily De la señal cruda a un conjunto de datos tabular:\\ ETL, calidad, selección de instancias y de características,\\ visualización y modelo de referencia\par}
}

\vfill
{\sffamily
\begin{tabular}{@{}l@{\hspace{1.2em}}l@{}}
  \color{acento}\bfseries Autor  & César Capelo \\[2pt]
  \color{acento}\bfseries Titulación & Grado en Ingeniería Informática en Tecnologías de la Información \\[2pt]
  \color{acento}\bfseries Profesor & Luciano Sánchez Ramos \\[2pt]
  \color{acento}\bfseries Tutor del TFG & Nahuel Costa Cortez \\[2pt]
  \color{acento}\bfseries Datos & Brugada-HUCA (PhysioNet, v1.0.0) \\[2pt]
  \color{acento}\bfseries Código & \url{https://github.com/CesarCapelo/BrugadaInteligenciaDeNegocio} \\[2pt]
  \color{acento}\bfseries Fecha & Noviembre de 2026 \\
\end{tabular}\par}
\end{titlepage}

% ---------------------------------------------------------------------
\chapter*{Resumen}
\addcontentsline{toc}{chapter}{Resumen}

El síndrome de Brugada es una canalopatía hereditaria asociada a arritmias ventriculares y muerte súbita en personas jóvenes con un corazón estructuralmente sano. Su diagnóstico depende de un patrón del electrocardiograma en las derivaciones precordiales derechas, el cual es dinámico, intermitente y, en la mayoría de los casos, se presenta de una forma sutil. Antes de entrenar cualquier modelo que ayude a realizar un diagnóstico, cada electrocardiograma, con sus 12 derivaciones, ha de convertirse a un conjunto pequeño de variables que sean fiables y que tengan un significado clínico. Dentro de este trabajo vamos a tratar de abordar la fase de preparación de los datos, dentro del proceso de descubrimiento de conocimiento, el KDD.

A partir del conjunto de datos Brugada-HUCA, el cual cuenta con un registro de 363 pacientes de los cuales 76 de ellos son positivos, con una duración de electrocardiograma de 12 segundos y muestreado a \SI{100}{\hertz}, se plantea construir el proceso completo y reproducible desde que se recibe la señal cruda hasta llegar a tener un conjunto tabular de $363 \times 72$. El proceso consta de seis fases:

\begin{enumerate}
  \item Un proceso ETL con control de calidad que detecta si hay un registro duplicado en el índice publicado y caracteriza la saturación, el ruido y la deriva que tiene cada registro.
  \item Un preprocesado de la señal con filtrado de fase cero, detección de latidos, segmentación y alineamiento.
  \item Una selección de instancias a nivel de latido y la construcción de un latido medio robusto por paciente.
  \item La extracción de las 72 características clínicas y la comparación de cuatro métodos de selección de características.
  \item La exploración de la estructura de los datos con PCA, t-SNE y UMAP.
  \item Un modelo de clasificación lineal de referencia.
\end{enumerate}

Además, todo lo que ajustamos con los datos se hace dentro de una validación cruzada anidada a nivel de paciente, de forma que no exista ningún tipo de fuga de información hacia el conjunto de test.

El modelo de referencia que alcanzo a realizar llega a tener un AUROC de \num{0.928} (IC 95\,\%: \num{0.88} a \num{0.96}), con una sensibilidad de \num{0.87} y especificidad \num{0.88}, frente al \num{0.75} que obtenemos de una regla que está basada en criterios clínicos. Obtuve también que la variable más discriminante de todas es la pendiente del segmento ST, no su elevación absoluta (AUC \num{0.84} frente a \num{0.53}). Además, con solo cinco variables, todas ellas del patrón de Brugada en V1--V2, somos capaces de conservar un \SI{99}{\percent} del rendimiento que sacamos con el modelo completo. Los cuatro métodos de selección rinden prácticamente igual, por lo que basta con uno sencillo y barato, como un filtro ANOVA. La selección de instancias descarta solo el \SI{0.9}{\percent} de los latidos, sin llegar a alterar el resultado, y ninguna proyección en dos dimensiones es capaz de separar las clases. Esta primera entrega de teoría es una parte de lo que constituye la memoria completa de mi futuro TFG cuando esté completo.

\tableofcontents
\listoffigures
\listoftables
\cleardoublepage
\pagenumbering{arabic}

% ---------------------------------------------------------------------
\clearpage
\section{Introducción}\label{cap:introduccion}

\subsection{Motivación}

El síndrome de Brugada es una enfermedad eléctrica del corazón de origen genético, la cual fue descrita por los hermanos Brugada en el año 1992, la cual predispone a la fibrilación ventricular y a la muerte súbita en personas jóvenes sin que estas cuenten con alguna cardiopatía estructural. En muchas ocasiones, la primera manifestación es el propio evento arrítmico. Por eso, es muy importante identificar el patrón electrocardiográfico en los ECG rutinarios, para que de esta forma los médicos sean capaces de estudiar al paciente y a su familia, evitar los fármacos o situaciones que puedan llegar a desencadenar arritmias, e incluso llegar a indicar un desfibrilador implantable.

El diagnóstico para esta cardiopatía se basa en un patrón característico en las derivaciones precordiales derechas V1 y V2, caracterizado por una elevación del segmento ST con un descenso convexo, seguido de una onda T negativa. Sin embargo, este patrón no siempre es el mismo: cuando aparece puede llegar a ser bastante sutil y puede llegar a confundirse con algunas variantes benignas como puede llegar a ser el bloqueo incompleto de la rama derecha. En el conjunto de datos que usaré dentro de este trabajo, de los 76 pacientes de la clase positiva, solo 23 de ellos tienen marcado el ECG basal como patológico. Es debido a esta gran dificultad para su diagnóstico que el uso del aprendizaje automático puede llegar a ser muy interesante.

La mayoría de los trabajos publicados acerca de esta enfermedad se centran en el modelo, pero en este problema, como en la mayoría de los problemas reales que nos encontramos dentro del campo de la Inteligencia de Negocio, el hecho de obtener o no un buen resultado depende de la calidad de los datos que vayamos a utilizar. Por eso, partimos de un ECG de 12 derivaciones de \SI{12}{\second} a \SI{100}{\hertz}, el cual contiene \num{14400} valores, que tienen ruido, pueden tener deriva en la línea de base, latidos anómalos, y el patrón que lo define dura apenas \SI{200}{\milli\second}, para tratar de convertirlo en unas pocas variables que sean fiables y que tengan de verdad un significado clínico. Es esta parte del KDD la que más esfuerzo exige y la que de verdad condiciona que cualquier modelo que entrenemos más tarde sea capaz de aprender.

\subsection{Objetivos}\label{sec:objetivos}

El objetivo general es el de construir y validar un proceso completo y reproducible capaz de transformar las señales ECG crudas del \emph{dataset} Brugada-HUCA en un conjunto de datos tabular listo para poder ser analizado. Los objetivos se corresponden con las siguientes fases del proceso:

\begin{enumerate}[label=\textbf{O\arabic*}, leftmargin=2.4em]
  \item \textbf{Proceso ETL y control de calidad}: extraer los registros, integrar estos con sus metadatos y detectar registros que estén duplicados, que tengan saturación, ruido o deriva.
  \item \textbf{Preprocesado de la señal}: filtrado, detección de latidos, segmentación y alineamiento, comparando varias alternativas para cada paso.
  \item \textbf{Selección de instancias}: rechazar los latidos que sean anómalos y construir un latido medio robusto por paciente.
  \item \textbf{Extracción y selección de características}: definir variables clínicas sobre el latido medio y comparar métodos de selección de características por rendimiento y coste.
  \item \textbf{Visualización}: explorar la propia estructura de los datos, usando técnicas de reducción de dimensionalidad, y cuantificar lo que están mostrando.
  \item \textbf{Modelo de referencia}: medir el valor que de verdad proporcionan estos datos, entrenando un clasificador sencillo, usando un protocolo de evaluación en el que no se produzca ninguna fuga de datos.
\end{enumerate}

\subsection{Contribuciones}

Además de cubrir los objetivos que acabo de nombrar, el trabajo también aporta:

\begin{itemize}
  \item Un proceso reproducible, que se encuentra en un repositorio abierto, con un comando por análisis, semillas fijas y una huella criptográfica del conjunto de datos.
  \item Un protocolo de evaluación sin fuga de información: todo lo que se ajusta con datos (normalización, selección de variables, hiperparámetros y umbral de decisión) se calcula dentro de cada partición de entrenamiento.
  \item Una comparación de métodos de selección de características de dos familias distintas (filtros y embebidos), evaluados sin sesgo de selección.
  \item Una confirmación cuantitativa de que la forma del segmento ST (su pendiente) discrimina mucho mejor que su elevación absoluta. No es una idea nueva, ya que coincide con criterios clínicos como el índice de Corrado, pero aquí se confirma con datos públicos y con métodos que no partían de ese conocimiento.
\end{itemize}

% ---------------------------------------------------------------------
\chapter{Fundamentos: el electrocardiograma y el síndrome de Brugada}\label{cap:fundamentos}

Este capítulo resume los conceptos clínicos necesarios para entender las decisiones de los capítulos siguientes: qué mide un ECG, en qué consiste el síndrome de Brugada y por qué su patrón es difícil de detectar.

\section{El electrocardiograma}

El electrocardiograma (ECG) registra en la superficie del cuerpo la actividad eléctrica que acompaña a la contracción del corazón. Cada latido produce una secuencia de ondas (\cref{fig:latido}):

\begin{itemize}
  \item la \textbf{onda P}, despolarización de las aurículas;
  \item el \textbf{complejo QRS}, despolarización de los ventrículos, con una duración normal por debajo de unos \SI{120}{\milli\second};
  \item el \textbf{punto J}, que marca el final del QRS y el comienzo del segmento ST;
  \item el \textbf{segmento ST}, una meseta normalmente isoeléctrica durante la cual los ventrículos permanecen despolarizados;
  \item la \textbf{onda T}, repolarización de los ventrículos.
\end{itemize}

El ECG estándar tiene 12 derivaciones, cada una de las cuales observa el mismo fenómeno eléctrico desde un ángulo distinto: seis de miembros (I, II, III, aVR, aVL, aVF) en el plano frontal y seis precordiales (V1--V6) en el plano horizontal. Las \textbf{precordiales derechas V1 y V2} se sitúan sobre el tracto de salida del ventrículo derecho, que es la región donde se origina el síndrome de Brugada. Por eso son las derivaciones en las que aparece su patrón y en las que se centra este trabajo.

\begin{figure}[htbp]
\centering
\begin{tikzpicture}[x=1.15cm, y=1.9cm, font=\sffamily\footnotesize]
  % rejilla
  \foreach \x in {0,0.4,...,10.01} \draw[acentoclaro, line width=0.3pt] (\x,-0.55) -- (\x,1.35);
  \foreach \y in {-0.5,-0.25,...,1.3} \draw[acentoclaro, line width=0.3pt] (0,\y) -- (10,\y);
  \draw[gristexto!50, dashed, line width=0.4pt] (0,0) -- (10,0);
  % latido normal
  \draw[acento, line width=1.3pt, line join=round]
    (0,0) -- (0.6,0) .. controls (0.9,0.2) and (1.3,0.2) .. (1.6,0) -- (2.6,0)
    -- (2.8,-0.12) -- (3.15,1.15) -- (3.5,-0.38) -- (3.72,0) -- (4.7,0.02)
    .. controls (5.3,0.05) and (5.6,0.48) .. (6.1,0.42) .. controls (6.5,0.38) and (6.8,0) .. (7.3,0) -- (10,0);
  % etiquetas de ondas
  \node[above] at (1.1,0.16) {P};
  \node[above] at (3.15,1.15) {R};
  \node[below left] at (2.85,-0.12) {Q};
  \node[below right] at (3.45,-0.38) {S};
  \node[above] at (6.1,0.45) {T};
  \fill[brugada] (3.72,0) circle (1.6pt) node[below right, text=brugada] {J};
  % llaves
  \draw[decorate, decoration={brace, amplitude=4pt, mirror}, gristexto] (2.6,-0.62) -- (3.72,-0.62) node[midway, below=4pt] {QRS};
  \draw[decorate, decoration={brace, amplitude=4pt, mirror}, gristexto] (3.72,-0.62) -- (5.0,-0.62) node[midway, below=4pt] {segmento ST};
  \draw[decorate, decoration={brace, amplitude=4pt}, gristexto] (0.6,0.42) -- (2.6,0.42) node[midway, above=4pt] {intervalo PR};
  % ventana de análisis
  \draw[brugada!70, line width=0.5pt, dashed] (3.72,-0.5) -- (3.72,1.3);
  \draw[brugada!70, line width=0.5pt, dashed] (4.62,-0.5) -- (4.62,1.3);
  \node[text=brugada, font=\sffamily\scriptsize, align=center] at (4.17,1.12) {J a J+80\\ ms};
\end{tikzpicture}
\caption[Anatomía de un latido en el ECG]{Anatomía esquemática de un latido. Cada cuadrícula horizontal representa \SI{40}{\milli\second}. El punto J (en rojo) separa la despolarización ventricular (QRS) de la repolarización (segmento ST y onda T). La ventana discontinua, de J a J+\SI{80}{\milli\second}, es la región donde se manifiesta el patrón de Brugada y donde se miden muchas de las características del \cref{cap:caracteristicas}.}
\label{fig:latido}
\end{figure}

\section{El síndrome de Brugada}

\subsection{Mecanismo, epidemiología y presentación clínica}

El síndrome de Brugada (SBr) es una \textbf{canalopatía}: una enfermedad de los canales iónicos de la membrana de las células cardíacas, sin alteración estructural visible del corazón. En torno a una quinta parte de los pacientes presenta variantes con pérdida de función en el gen \emph{SCN5A}, que codifica el canal de sodio cardíaco. La reducción de la corriente de sodio se expresa sobre todo en el epicardio del tracto de salida del ventrículo derecho \cite{brugada2018jacc,antzelevitch2016}. Coexisten dos explicaciones, no excluyentes, de cómo se genera el patrón: un gradiente de repolarización entre el epicardio y el endocardio de esa región (hipótesis de repolarización) y un retraso de conducción localizado (hipótesis de despolarización).

Su prevalencia se estima entre 1 de cada \num{2000} y 1 de cada \num{5000} personas, con valores más altos en el sudeste asiático. Es entre ocho y diez veces más frecuente en varones, y suele diagnosticarse en torno a los 40~años \cite{brugada2018jacc}. Muchos pacientes están asintomáticos. Otros debutan con síncope o con fibrilación ventricular, típicamente en reposo, durante el sueño o con fiebre.

\subsection{Patrones electrocardiográficos}

Los consensos internacionales describen dos patrones en V1--V3 \cite{wilde2002,bayesdeluna2012,priori2013} (\cref{fig:patrones}):

\begin{itemize}
  \item \textbf{Tipo 1 o \eng{coved} (convexo)}: elevación del punto J de al menos \SI{2}{\milli\metre} (\SI{0.2}{\milli\volt}), seguida de un segmento ST descendente y convexo que termina en una \textbf{onda T negativa} y simétrica. Es el único patrón diagnóstico.
  \item \textbf{Tipo 2 o \eng{saddleback} (en silla de montar)}: onda r' elevada, seguida de un segmento ST cóncavo que se mantiene elevado al menos \SI{0.5}{\milli\metre} y de una onda T positiva o aplanada. Es sugestivo, pero no diagnóstico. El antiguo tipo 3 se integró en él en 2012 \cite{bayesdeluna2012}.
\end{itemize}

Para distinguir el tipo 2 de variantes benignas con onda r' (bloqueo incompleto de rama derecha, deportistas) se han propuesto criterios cuantitativos sobre la forma del complejo: el ángulo $\beta$ entre las dos ramas de la r' \cite{chevallier2011} y la anchura de la base del triángulo que forma la r' \cite{serra2014}. Para distinguir el patrón tipo 1 de la repolarización precoz de los deportistas se propuso el \textbf{índice de Corrado}, el cociente entre la elevación del ST en el punto J y a \SI{80}{\milli\second} ($\mathrm{ST_J}/\mathrm{ST_{80}}$): un valor mayor que 1 indica un ST descendente, propio del patrón de Brugada \cite{zorzi2015}. Todos estos criterios miden la \emph{forma} de la transición entre el QRS y el ST, no su altura. Los datos de este trabajo confirman empíricamente esa idea (\cref{cap:caracteristicas}).

\begin{figure}[htbp]
\centering
\begin{tikzpicture}[x=0.95cm, y=1.55cm, font=\sffamily\footnotesize]
  \foreach \ox/\tit in {0/{Normal (V1)}, 5.2/{Tipo 1 (\eng{coved})}, 10.4/{Tipo 2 (\eng{saddleback})}} {
    \begin{scope}[shift={(\ox,0)}]
      \foreach \x in {0,0.4,...,4.41} \draw[acentoclaro, line width=0.3pt] (\x,-1.1) -- (\x,1.3);
      \foreach \y in {-1.0,-0.75,...,1.3} \draw[acentoclaro, line width=0.3pt] (0,\y) -- (4.4,\y);
      \draw[gristexto!50, dashed, line width=0.4pt] (0,0) -- (4.4,0);
      \node[anchor=south west, font=\sffamily\small\bfseries, text=acento] at (0,1.32) {\tit};
    \end{scope}}
  % normal en V1: rS, ST isoeléctrico, T positiva baja
  \draw[normal, line width=1.3pt, line join=round] (0,0) -- (0.9,0) -- (1.05,0.25) -- (1.25,-0.95) -- (1.45,0) -- (2.1,0.03)
    .. controls (2.5,0.08) and (2.7,0.3) .. (3.0,0.28) .. controls (3.3,0.26) and (3.5,0) .. (3.8,0) -- (4.4,0);
  % tipo 1: r' alta que desciende en rampa convexa hasta T negativa
  \begin{scope}[shift={(5.2,0)}]
    \draw[brugada, line width=1.3pt, line join=round] (0,0) -- (0.9,0) -- (1.05,0.2) -- (1.22,-0.7) -- (1.42,0.95)
      .. controls (1.9,0.92) and (2.35,0.65) .. (2.65,0.05) .. controls (2.85,-0.35) and (3.1,-0.5) .. (3.35,-0.3)
      .. controls (3.55,-0.12) and (3.7,0) .. (3.95,0) -- (4.4,0);
    \draw[gristexto, -{Latex[length=1.6mm]}, line width=0.5pt] (2.75,1.08) node[right, font=\sffamily\scriptsize, text=gristexto] {ST descendente} -- (2.12,0.84);
    \draw[gristexto, -{Latex[length=1.6mm]}, line width=0.5pt] (3.6,-0.85) node[below, font=\sffamily\scriptsize, text=gristexto] {T negativa} -- (3.2,-0.47);
    \draw[brugada, line width=0.5pt] (1.42,0) -- (1.42,0.95); \node[left, font=\sffamily\scriptsize, text=brugada, fill=white, inner sep=1pt] at (1.36,0.62) {$\geq$\,2 mm};
  \end{scope}
  % tipo 2: r' y ST en silla de montar, T positiva
  \begin{scope}[shift={(10.4,0)}]
    \draw[acento, line width=1.3pt, line join=round] (0,0) -- (0.9,0) -- (1.05,0.2) -- (1.22,-0.7) -- (1.42,0.85)
      .. controls (1.6,0.35) and (1.8,0.22) .. (2.1,0.24) .. controls (2.5,0.27) and (2.7,0.6) .. (3.0,0.58)
      .. controls (3.3,0.55) and (3.5,0) .. (3.85,0) -- (4.4,0);
    \draw[gristexto, -{Latex[length=1.6mm]}, line width=0.5pt] (2.45,-0.72) node[right, font=\sffamily\scriptsize, text=gristexto] {ST cóncavo elevado} -- (2.07,0.17);
  \end{scope}
\end{tikzpicture}
\caption[Patrones de Brugada en V1--V2]{Representación esquemática de los patrones de Brugada en V1--V2 frente a un complejo normal de V1. El patrón tipo 1 combina una elevación del punto J de al menos \SI{0.2}{\milli\volt}, un ST descendente y convexo y una onda T negativa. El tipo 2 presenta una r' y un ST cóncavo en silla de montar con T positiva. Adaptado de los criterios de consenso \cite{bayesdeluna2012,priori2013}.}
\label{fig:patrones}
\end{figure}

\subsection{Diagnóstico}

El consenso HRS/EHRA/APHRS de 2013 establece que el SBr se diagnostica cuando se observa un patrón tipo~1 en al menos una derivación precordial derecha (V1 o V2, en posición estándar o elevada hasta el segundo espacio intercostal), ya sea espontáneo o inducido por un bloqueante de los canales de sodio como la ajmalina o la flecainida \cite{priori2013}. El consenso de 2016 propuso además una puntuación (\eng{Shanghai score}) que combina el ECG, la clínica, los antecedentes familiares y la genética \cite{antzelevitch2016}. Las guías europeas de 2022 restringieron el valor del patrón inducido por fármacos, que por sí solo ya no basta para el diagnóstico \cite{zeppenfeld2022}.

\section{Implicaciones para el análisis de datos}

\begin{table}[!b]
\centering
\caption[Traducción de los criterios clínicos a variables]{Traducción de los criterios clínicos del patrón de Brugada a variables calculables sobre el latido medio. La definición precisa de cada variable está en el \cref{ap:catalogo}.}
\label{tab:criterios}
\small
\begin{tabularx}{\linewidth}{@{}p{4.2cm}Lp{3.0cm}@{}}
\toprule
\textbf{Criterio clínico} & \textbf{Variable en este trabajo} & \textbf{Nombre} \\
\midrule
Elevación del punto J $\geq$ \SI{0.2}{\milli\volt} & Amplitud del latido en el punto J y a \SIlist{40;80}{\milli\second} & \var{st_j}, \var{st_j40}, \var{st_j80} \\
ST descendente y convexo (\eng{coved}); índice de Corrado & Pendiente del ST entre J y J+\SI{80}{\milli\second}; pendiente final del QRS & \var{st_slope}, \var{qrs_end_slope} \\
Onda T negativa & Mínimo de la onda T e indicador de negatividad & \var{t_neg}, \var{t_negative} \\
Onda r' prominente & Máximo positivo tras la onda S & \var{rprime_amp} \\
Ensanchamiento del QRS & Duración del QRS a partir de la energía de la derivada & \var{qrs_dur_ms} \\
Elevación relativa del ST & Cociente entre el ST y la amplitud del QRS & \var{st_r_ratio} \\
\bottomrule
\end{tabularx}
\end{table}

Las características clínicas del síndrome tienen consecuencias directas sobre cómo hay que preparar los datos:

\begin{enumerate}
  \item \textbf{El patrón es local en el espacio y en el tiempo.} Está en dos derivaciones (V1--V2) y ocupa una ventana de unos \SI{200}{\milli\second} alrededor del punto J. La segmentación y el alineamiento de latidos deben ser precisos en esa zona.
  \item \textbf{El patrón es intermitente y a menudo incompleto.} La mayoría de los pacientes del conjunto de datos no muestra el tipo 1 en su ECG basal, así que las variables deben captar tendencias sutiles y no solo comprobar si se cumple un umbral clínico.
  \item \textbf{Los criterios clínicos son medibles.} Cada elemento del patrón se puede traducir en una variable calculable sobre el latido (\cref{tab:criterios}). Esto permite construir características interpretables y comparar un sistema de reglas explícitas con modelos aprendidos.
\end{enumerate}


% ---------------------------------------------------------------------
\chapter{Datos y proceso ETL}\label{cap:datos}

Dentro de este capítulo vamos a explorar el conjunto de datos, el proceso de extracción, transformación y carga, el control de calidad y el protocolo de evaluación que vamos a seguir durante todo el trabajo. En la \cref{fig:pipeline} se resume el proceso completo:

\begin{figure}[htbp]
\centering
\begin{tikzpicture}[
  font=\sffamily\scriptsize, node distance=3.2mm and 3mm,
  paso/.style={draw=acento, fill=acentoclaro!55, rounded corners=3pt, text width=24mm, minimum height=13mm, align=center, line width=0.6pt},
  dato/.style={draw=gristexto!80, fill=white, rounded corners=6pt, text width=24mm, minimum height=9mm, align=center},
  etq/.style={font=\sffamily\tiny\bfseries, text=acento, anchor=south east},
  f/.style={-{Latex[length=2mm]}, acento, line width=0.7pt}]
  \node[dato] (raw) {\textbf{363 ECG crudos}\\ 12 der.\ $\times$ 1200 muestras};
  \node[paso, right=of raw] (etl) {\textbf{ETL y calidad}\\ duplicados, saturación, ruido, deriva};
  \node[paso, right=of etl] (fil) {\textbf{Filtrado}\\ paso banda \SIrange{0.5}{40}{\hertz}, fase cero};
  \node[paso, right=of fil] (det) {\textbf{Detección de latidos}\\ picos R en II (NeuroKit2)};
  \node[paso, below=8mm of det] (seg) {\textbf{Segmentación y alineamiento}\\ \SIrange{-250}{450}{\milli\second}, $\pm$\SI{30}{\milli\second}};
  \node[paso, left=of seg] (sel) {\textbf{Selección de instancias}\\ rechazo de latidos anómalos};
  \node[paso, left=of sel] (med) {\textbf{Latido medio}\\ mediana por paciente, $70\times12$};
  \node[paso, left=of med] (car) {\textbf{Extracción de características}\\ 72 variables clínicas};
  \node[dato, below=8mm of car] (tab) {\textbf{Tabla} $\mathbf{363 \times 72}$};
  \node[paso, right=of tab, fill=grisclaro, draw=gristexto] (fs) {\textbf{Selección de características}};
  \node[paso, right=of fs, fill=grisclaro, draw=gristexto] (vis) {\textbf{Visualización}\\ PCA, t-SNE, UMAP};
  \node[paso, right=of vis, fill=grisclaro, draw=gristexto] (mod) {\textbf{Modelo de referencia}\\ LDA, CV anidada};
  \draw[f] (raw) -- (etl); \draw[f] (etl) -- (fil); \draw[f] (fil) -- (det); \draw[f] (det) -- (seg);
  \draw[f] (seg) -- (sel); \draw[f] (sel) -- (med); \draw[f] (med) -- (car); \draw[f] (car) -- (tab);
  \draw[f] (tab) -- (fs); \draw[f] (fs) -- (vis); \draw[f] (vis) -- (mod);
  \node[etq] at (etl.north east) {Cap.\ \ref{cap:datos}}; \node[etq] at (fil.north east) {Cap.\ \ref{cap:preprocesado}}; \node[etq] at (det.north east) {Cap.\ \ref{cap:preprocesado}};
  \node[etq] at (seg.north east) {Cap.\ \ref{cap:preprocesado}}; \node[etq] at (sel.north east) {Cap.\ \ref{cap:instancias}}; \node[etq] at (med.north east) {Cap.\ \ref{cap:instancias}};
  \node[etq] at (car.north east) {Cap.\ \ref{cap:caracteristicas}}; \node[etq] at (fs.north east) {Cap.\ \ref{cap:caracteristicas}}; \node[etq] at (vis.north east) {Cap.\ \ref{cap:visualizacion}}; \node[etq] at (mod.north east) {Cap.\ \ref{cap:modelo}};
\end{tikzpicture}
\caption[Visión general del proceso]{Visión general del proceso, desde la señal cruda hasta la tabla de características y su análisis. En azul, las etapas que construyen el conjunto de datos. En gris, las que lo analizan. Cada etapa indica el capítulo en el que se desarrolla.}
\label{fig:pipeline}
\end{figure}

\section{El conjunto de datos Brugada-HUCA}

Brugada-HUCA es un conjunto de datos público, publicado en PhysioNet \cite{costacortez2026huca,goldberger2000} con la licencia CC BY-SA~4.0. Fue recogido por el Servicio de Cardiología del Hospital Universitario Central de Asturias, y las etiquetas fueron asignadas por electrofisiólogos siguiendo los criterios de consenso, con revisión conjunta para los casos que eran dudosos. Este contiene un ECG de 12 derivaciones por paciente, de \SI{12}{\second} de duración y muestreado a \SI{100}{\hertz}, en formato WFDB. Y luego tenemos el fichero \cod{metadata.csv}, que nos da tres variables por paciente (\cref{tab:composicion}):

\begin{itemize}
  \item \var{brugada}: que nos indica el diagnóstico para este paciente, 0 si está sano, 1 si tiene el síndrome de Brugada y 2 si se trata de un caso atípico.
  \item \var{basal_pattern}: si el ECG basal se considera patológico.
  \item \var{sudden_death}: si hubo algún antecedente de muerte súbita o de evento arrítmico.
\end{itemize}

\begin{table}[htbp]
\centering
\caption[Composición del conjunto de datos]{Composición del conjunto de datos Brugada-HUCA.}
\label{tab:composicion}
\small\setlength{\tabcolsep}{5pt}
\begin{tabular}{@{}lS[table-format=3.0]S[table-format=2.0]S[table-format=2.0]S[table-format=3.1]@{}}
\toprule
\textbf{Diagnóstico} & {\textbf{Pacientes}} & {\textbf{Basal patológico}} & {\textbf{Muerte súbita}} & {\textbf{\% del total}} \\
\midrule
0 --- Normal & 287 & 23 & 6 & 79.1 \\
1 --- Brugada confirmado & 69 & 19 & 4 & 19.0 \\
2 --- Atípico & 7 & 4 & 1 & 1.9 \\
\midrule
Total & 363 & 46 & 11 & 100.0 \\
\bottomrule
\end{tabular}
\end{table}

\section{Proceso ETL}

\paragraph{Extracción.} Tenemos un \eng{script} que se descarga los ficheros de PhysioNet y los lee con la librería \cod{wfdb}. Luego, para cada registro, el cargador comprueba que las 12 derivaciones están en orden (I, II, III, aVR, aVL, aVF, V1--V6) y que además la señal tenga exactamente $1200 \times 12$ muestras. Si alguna de las comprobaciones que se van haciendo llegase a fallar, el proceso se detendría para no seguir con datos incorrectos.

\paragraph{Transformación.} Aquí los metadatos se integran con la señal a través del identificador propio que tiene cada paciente y construimos la etiqueta binaria del problema: positivo si \var{brugada} es 1 o 2, y negativo si \var{brugada} es 0. Lo que termina dándonos 76 pacientes positivos, que representan un \SI{20.9}{\percent} del total, y 287 negativos. Como se ve, los 7 casos atípicos que hay los incluyo como positivos; esto lo explicaré más adelante (\cref{sec:etiqueta}) y analizaré su efecto.

\paragraph{Carga.} Todas las señales se guardan dentro de un único \eng{array} de $363 \times 1200 \times 12$ en formato comprimido; así nos libramos de tener que volver a leer los ficheros WFDB. Además, calculamos una huella SHA-256 de todos los ficheros de señal que identifica en qué versión exacta se encuentran los datos. También las tablas intermedias (los latidos medios, las características y los indicadores de calidad por paciente) se guardan en ficheros CSV y NumPy versionados junto al código.

\section{Control de calidad}\label{sec:calidad}

En la \cref{tab:calidad} resumo las comprobaciones realizadas, y en la \cref{fig:calidad} se muestra la composición del conjunto, el espectro medio de la señal y la distribución de la frecuencia cardíaca.

\begin{table}[htbp]
\centering
\caption[Comprobaciones de calidad de los datos]{Comprobaciones de calidad de los datos y acciones tomadas.}
\label{tab:calidad}
\small
\begin{tabularx}{\linewidth}{@{}p{3.3cm}Lp{3.9cm}@{}}
\toprule
\textbf{Comprobación} & \textbf{Resultado} & \textbf{Acción} \\
\midrule
Registros duplicados & El índice \cod{RECORDS} publicado contiene 364 entradas: el paciente 596382 aparece dos veces & Se elimina la entrada duplicada (363 registros únicos) \\
Integridad referencial & Los 363 identificadores coinciden entre señales y metadatos & Ninguna \\
Valores ausentes y derivaciones planas & Ninguno & Ninguna \\
Saturación ($|x| \geq \SI{3.2}{\milli\volt}$) & 11 registros en alguna derivación. En V1--V2, 3 pacientes con entre 2 y 17 muestras, casi todas artefactos (picos aislados de signo alterno o al inicio del registro); solo 2 muestras coinciden con un QRS & Se conservan: sin esos 3 pacientes el AUROC de referencia no cambia (\num{0.929} frente a \num{0.928}) \\
Ruido de alta frecuencia & Mediana del \SI{0.02}{\percent} de la potencia de V1 por encima de \SI{40}{\hertz}; percentil 95 del \SI{0.51}{\percent} & Filtro paso bajo a \SI{40}{\hertz} \\
Deriva de la línea de base & Mediana del \SI{0.56}{\percent} de la potencia de V1 por debajo de \SI{0.5}{\hertz} & Filtro paso alto a \SI{0.5}{\hertz} \\
Calidad de los latidos & Mediana de 13 latidos válidos por registro; correlación media con la plantilla de \num{0.98} & Los 363 registros superan el control \\
Frecuencia cardíaca & Mediana de 69 lpm (de 40 a 157); sin diferencias entre clases (Mann-Whitney, $p = \num{0.11}$) & Ninguna \\
\bottomrule
\end{tabularx}
\end{table}

\begin{figure}[htbp]
\centering
\includegraphics[width=\linewidth]{calidad_datos}
\caption[Composición, espectro y ritmo]{(a) Pacientes por diagnóstico; la parte rayada corresponde al ECG basal patológico. (b) Densidad espectral de potencia media de V1 y V2 (método de Welch): la energía se concentra por debajo de \SI{30}{\hertz} y cae unos tres órdenes de magnitud hasta \SI{40}{\hertz}. La banda gris indica el ruido que elimina el filtro paso bajo. (c) La frecuencia cardíaca no difiere entre clases, por lo que no puede actuar como variable de confusión.}
\label{fig:calidad}
\end{figure}

\begin{hallazgo}
Como vemos, el conjunto de datos tiene una buena calidad: no hay valores ausentes, el ruido y la deriva que tenemos son pequeños, y todos los registros nos dejan hacerles el latido medio fiable. El único defecto que tiene es un paciente duplicado, el cual se eliminó, y que de no haberlo hecho podría haber quedado en el entrenamiento y en el test también, introduciendo de esta manera una fuga de información.
\end{hallazgo}

\section{Definición de la etiqueta}\label{sec:etiqueta}

Los 7 casos atípicos que tenemos dentro del dataset corresponden a pacientes los cuales están siendo estudiados por sospecha de Brugada, pero que no llegan a cumplir con todos los criterios de confirmación. Por mi parte, los considero positivos, ya que desde un punto de vista clínico están más cerca de parecerse a alguien con Brugada que a alguien sano, desde el punto de vista del procedimiento que hay que seguir, ya que estos casos son los que merecen un estudio adicional como los positivos en Brugada, y por eso deberían de estar juntos al desencadenar el mismo resultado. Comprobé que esta decisión no condicionaba las conclusiones, y evalué dos alternativas con la regresión logística, que con la etiqueta elegida obtiene un AUROC de \num{0.919} (\cref{cap:modelo}). En la primera decidí excluir a los atípicos, dándonos un AUROC de \num{0.923} ($\pm$\num{0.045}), y en cambio, si los trato como negativos, se llega a reducir hasta el \num{0.903} ($\pm$\num{0.042}). Como se ve en los datos obtenidos, el descenso nos indica que estos casos atípicos se parecen más a los pacientes con Brugada que a los que no, lo que apoya el hecho de incluirlos dentro de la clase positiva.

\section{Protocolo de evaluación común}\label{sec:protocolo}

Ya que varias de las decisiones que se van a tomar más adelante se van a hacer sobre un clasificador, es imperativo que todas sigan el mismo protocolo, para hacer que las comparaciones sean justas y no estén infladas. El protocolo es el siguiente:

\begin{enumerate}
  \item \textbf{Particiones a nivel de paciente:} usamos una validación cruzada estratificada de 5 particiones, repetida 3 veces usando una semilla fija, lo que nos da un total de 15 \eng{folds}. Además, como tenemos solo un ECG por paciente, nos aseguramos de que cada paciente aparece exactamente una sola vez en el conjunto de test de cada repetición.
  \item \textbf{Validación anidada:} todo lo que el modelo aprende con los datos se ajusta solo con la partición de entrenamiento de cada \eng{fold}. De esta forma, lo que evitamos es el sesgo optimista que aparece cuando resulta que el modelo usa los mismos datos que con los que se evalúa \cite{varma2006,ambroise2002}.
  \item \textbf{Umbral de decisión:} este se elige maximizando el índice de Youden \cite{youden1950} sobre las predicciones de fuera de la muestra de la validación interna, es decir, el umbral que hace máxima la suma de la sensibilidad y la especificidad.
  \item \textbf{Desbalance:} como tenemos muchos más pacientes normales que con Brugada, un modelo que siempre diga que el paciente es normal acertaría un \SI{79}{\percent} de las veces sin llegar a detectar ningún caso. Para evitar esto, durante el entrenamiento cada clase recibe un peso que es inversamente proporcional a su frecuencia, de modo que equivocarse en un Brugada le cueste 3,8 veces más que equivocarse con uno normal. Además, uso métricas las cuales no dependen de la prevalencia, como lo son el AUROC y la exactitud balanceada, $\mathrm{BA} = (\mathrm{Sens} + \mathrm{Esp})/2$.
  \item \textbf{Incertidumbre:} se informa de la media $\pm$ la desviación típica del AUROC en los 15 \eng{folds}, y de un intervalo de confianza del \SI{95}{\percent} obtenido por \eng{bootstrap} sobre las predicciones que se encuentran fuera de la muestra.
\end{enumerate}

\begin{nota}[Un apunte]
Con 76 positivos, en cada \eng{fold} de test tenemos aproximadamente unos 15 pacientes con Brugada. La desviación típica del AUROC entre \eng{folds} está sobre el \num{0.05}, por lo que las diferencias que sean inferiores a \num{0.02} o \num{0.03} no las podemos interpretar directamente como mejoras reales sin antes haberlas pasado por un contraste estadístico.
\end{nota}

% ---------------------------------------------------------------------
\clearpage
\section{Preprocesado de la señal}\label{cap:preprocesado}

El preprocesado lo que hace es convertir cada registro en un conjunto de latidos limpios, los cuales estén sincronizados entre las derivaciones y alineados entre sí. Este consta de cuatro pasos deterministas, los cuales no usan ninguna etiqueta, y que por lo tanto no llegan a introducir ningún tipo de fuga de información, aunque se apliquen antes de particionar siquiera los datos. En la \cref{fig:senal} se muestra su efecto sobre el ECG de un paciente.

\figura[\linewidth]{senal_preprocesado}{Efecto del preprocesado sobre un paciente con síndrome de Brugada. (a) V2 cruda, con una deriva de la línea de base de \SI{0.69}{\milli\volt}. (b) La misma derivación tras el filtro paso banda de fase cero: la deriva desaparece y la morfología de cada latido se conserva. (c) Derivación II filtrada con los picos R detectados, que se usan para segmentar las 12 derivaciones.}{fig:senal}

\subsection{Filtrado}

Se aplica un filtro de Butterworth paso banda de orden 4 entre $f_1 = \SI{0.5}{\hertz}$ y $f_2 = \SI{40}{\hertz}$. Este filtro de Butterworth tiene una respuesta en amplitud máximamente plana en la banda de paso, por lo que no llega a alterar la amplitud relativa de las ondas que nos interesan. El paso alto lo que consigue es eliminar la deriva de la línea de base, y el paso bajo consigue eliminar el ruido muscular. Además, al estar recolectados estos datos aquí en Asturias, donde la frecuencia de la red eléctrica es de \SI{50}{\hertz}, esta coincide con la frecuencia de Nyquist del registro, por lo que la energía residual solo puede llegar a aparecer en los bordes de la banda, que es una zona la cual el paso bajo también atenúa.

El filtro lo aplicamos en dos direcciones, primero hacia delante y después hacia atrás. Las distorsiones de fase que se producen en cada pasada se cancelan, de modo que el filtro resultante tiene fase cero y una respuesta en amplitud igual al cuadrado de la original.
Esto de tener fase cero es esencial, ya que un filtro que retrasara la onda según su frecuencia lo que haría sería desplazar el punto J y el segmento ST, que es la parte que nos concierne, ya que es la que identifica al Brugada.

\subsection{Detección de latidos}

Los picos R se detectan usando el método por defecto de NeuroKit2 sobre la derivación II, la cual suele tener complejos QRS amplios y positivos. Se activa la corrección de artefactos, la cual elimina detecciones con intervalos RR fisiológicamente imposibles y recupera latidos perdidos. Si se detectan menos de 3 latidos, se repite este proceso sobre V5, V2 o I. Y ya por último, cada pico se reajusta al máximo absoluto de la señal dentro de una ventana de $\pm\SI{50}{\milli\second}$, con el fin de corregir pequeños desplazamientos que pueda haber sobre el detector.

\subsection{Segmentación y alineamiento}

Una vez localizados los picos R, recortamos cada latido en una ventana de \SI{250}{\milli\second} antes a \SI{450}{\milli\second} después del pico, lo que a \SI{100}{\hertz} son 70 muestras, y que incluye el final de la onda P, el QRS, el segmento ST y la onda T. El corte se hace en los mismos instantes en las 12 derivaciones, para que todas queden sincronizadas, y los latidos de los bordes que no caben enteros se descartan.

Como a \SI{100}{\hertz} un error de una o dos muestras en el pico R ya desplaza el latido 10 o \SI{20}{\milli\second}, y eso emborronaría el segmento ST al promediar, alineamos cada latido con la plantilla del paciente, esto sería la mediana de sus latidos. Entonces probamos a moverlo hasta 3 muestras hacia cada lado y nos quedamos con la posición en la que más se parece a la plantilla.

\subsection{Comparación de alternativas}\label{sec:ablacion-preprocesado}

Cada paso podría haberse hecho de otra forma; hay varias alternativas razonables. Para elegir entre una de ellas con datos, se aplicó el proceso que describimos completo a cada alternativa, y se midió el desempeño del clasificador de referencia: una regresión logística sobre las 72 características, evaluada con el protocolo de la \cref{sec:protocolo}. En la \cref{tab:ablacion} se ven los resultados:

\begin{table}[H]
\centering
\caption[Comparación de alternativas de preprocesado]{Comparación de alternativas en cada paso del preprocesado (AUROC medio $\pm$ desviación típica en 15 \emph{folds}). En negrita, la alternativa elegida.}
\label{tab:ablacion}
\small
\begin{tabular}{@{}l>{\raggedright\arraybackslash}p{6.6cm}S[table-format=1.3]S[table-format=1.3]@{}}
\toprule
\textbf{Paso} & \textbf{Alternativa} & {\textbf{AUROC}} & {\textbf{DT}} \\
\midrule
\multirow{4}{*}{Filtrado} & \textbf{Paso banda \SIrange{0.5}{40}{\hertz}} & 0.919 & 0.052 \\
 & Doble mediana (\SIlist{200;600}{\milli\second}) y paso bajo \SI{40}{\hertz} & 0.919 & 0.055 \\
 & Paso banda \SIrange{0.67}{30}{\hertz} & 0.914 & 0.055 \\
 & Sin filtrar & 0.913 & 0.056 \\
\midrule
\multirow{4}{*}{Detector de picos R} & \textbf{NeuroKit2} & 0.919 & 0.052 \\
 & Pan-Tompkins & 0.903 & 0.052 \\
 & Hamilton & 0.876 & 0.046 \\
 & Elgendi & 0.866 & 0.057 \\
\midrule
\multirow{3}{*}{Derivación de referencia} & \textbf{II} & 0.919 & 0.052 \\
 & V5 & 0.927 & 0.045 \\
 & V2 & 0.892 & 0.054 \\
\midrule
\multirow{2}{*}{Agregación de latidos} & \textbf{Mediana} & 0.919 & 0.052 \\
 & Media & 0.912 & 0.052 \\
\bottomrule
\end{tabular}
\end{table}

A la vista de estos resultados, podemos sacar tres conclusiones:

\begin{itemize}
  \item \textbf{El filtrado casi no influye}, ya que sin filtrarlo tenemos \num{0.913} y filtrándolo \num{0.919} de AUROC. Lo que nos indica que las señales con las que estamos trabajando ya son limpias y que el latido medio atenúa bastante el ruido residual.
  \item \textbf{El detector de latidos sí que influye}, dándonos hasta un \num{0.05} más de AUROC. Vemos que un pico que esté mal situado consigue desalinear los latidos, haciendo que el latido medio quede deformado en la zona QRS--ST, que es donde se encuentra el patrón de detección del Brugada. Vemos que NeuroKit2 combina un detector robusto con corrección de artefactos y es el que mejor funciona.
  \item \textbf{La derivación de referencia sí que importa.} Vemos que V2 es la peor referencia de todas, ya que su complejo QRS es el que se altera en pacientes con Brugada, lo que hace más difícil el localizar el pico. Vemos también que V5 saca un valor mayor que II, pero la diferencia que existe entre los \emph{folds} es mayor que lo que ganamos; por eso se mantiene II, que además es la referencia estándar en la práctica clínica y la elección que se decidió usar antes de ver estos resultados. Usar V5 en vez de II sería ajustar el proceso al conjunto de test, y esto es algo que queremos evitar a toda costa.
\end{itemize}

\begin{conclusion}
Al final, lo que buscamos con esto es precisión temporal, ya que, a la hora de la verdad, el Brugada ocupa solo unos pocos cientos de milisegundos alrededor del punto J, por lo que el hecho de que haya un poco de ruido o no, no llega a ser tan perjudicial como que los latidos no estén bien localizados o alineados.
\end{conclusion}

% ---------------------------------------------------------------------
\clearpage
\section{Selección de instancias y latido medio}\label{cap:instancias}

Tenemos que cada registro nos aporta entre 8 y 26 latidos completos, con una mediana de 13. Estos latidos son las \emph{instancias} a través de las cuales somos capaces de resumir a cada paciente. Antes de combinarlos, primero se eliminan los anómalos, como pueden ser latidos ectópicos o detecciones erróneas, en un proceso de edición no supervisada (\cref{sec:estado}), comparando cada latido con el prototipo de su propio paciente, sin usar la etiqueta de clase.

\subsection{Criterios de rechazo}

Se han definido dos criterios, los cuales se pueden combinar.

\paragraph{Criterio morfológico.} Se calcula la plantilla mediana de los latidos alineados de un registro en la derivación de referencia, y cada latido se conserva si su correlación de Pearson con la plantilla es igual o mayor que un umbral $\rho$, para el que se probaron los valores \num{0.5}; \num{0.6}; \num{0.7}; \num{0.8} y \num{0.9}. Detecta latidos con otra morfología, como las extrasístoles ventriculares, los artefactos de movimiento o los complejos mal detectados.

\paragraph{Criterio temporal.} Se descarta el latido si su intervalo RR previo se desvía más de un 20\,\% del RR mediano del registro. Detecta latidos prematuros y latidos tras una pausa, cuya morfología puede ser normal pero cuya repolarización está alterada.

En ambos casos, si el criterio dejara menos de 3 latidos se conservan todos (\emph{fallback}). En la práctica esto solo ocurre en algún paciente, sobre todo con los umbrales más exigentes.

\subsection{Construcción del latido medio}

Ya con los latidos conservados, se calcula, para cada derivación y cada instante, la \textbf{mediana} muestra a muestra. Después se resta el nivel isoeléctrico, estimado como la media del segmento PR (de \SIrange{-80}{-40}{\milli\second} respecto al pico R).

La mediana tiene un punto de ruptura del 50\,\%, por lo que aunque cerca de la mitad de los latidos fueran anómalos, esta seguiría representando a los normales. Sin embargo, esto no pasa con la media, que sí se ve arrastrada por un único latido extremo. El resultado que tenemos es una matriz de $70 \times 12$ por paciente, con una misma dimensión para todos, que es la base de las características. En la \cref{fig:segmentacion} vemos el proceso con un paciente del cual, de sus 14 latidos, 4 son anómalos.

\figura[\linewidth]{latidos_segmentacion}{Selección de instancias en un paciente (ID 3118448) con 4 latidos anómalos de 14. (a) Latidos segmentados en la derivación II, antes de alinear. (b) Tras el alineamiento, en rojo los descartados por el criterio morfológico ($\rho = 0{,}6$): morfologías claramente distintas de la plantilla. (c) Latidos conservados en V2 y latido medio resultante (en negro).}{fig:segmentacion}

\subsection{Resultados}

La \cref{fig:embudo} resume cuántos latidos se pierden en cada paso usando el criterio por defecto ($\rho = 0{,}6$), y en la \cref{tab:instancias} comparamos todos los criterios.

\begin{figure}[H]
\centering
\makebox[\linewidth][c]{\begin{tikzpicture}[font=\sffamily\small, x=1cm, y=1cm]
  \def\escala{0.00145}
  \foreach \n/\v/\y/\col/\lab in {1/4957/0/gristexto!35/Picos R detectados, 2/4776/-1.05/gristexto!35/Latidos completos, 3/4735/-2.1/acento!45/{Conservados ($\rho \geq 0{,}6$)}} {
    \fill[\col, rounded corners=2pt] (0,\y-0.36) rectangle ({\v*\escala},\y+0.36);
    \node[anchor=east] at (-0.2,\y) {\lab};
    \node[anchor=west, font=\sffamily\small\bfseries] at ({\v*\escala+0.15},\y) {\num{\v}};
  }
  \node[anchor=west, font=\sffamily\footnotesize, text=gristexto, align=left] at (8.15,-1.05) {$-181$ (\SI{3.7}{\percent})\\ bordes del registro};
  \node[anchor=west, font=\sffamily\footnotesize, text=gristexto, align=left] at (8.15,-2.1) {$-41$ (\SI{0.9}{\percent})\\ baja correlación};
  \draw[-{Latex[length=2mm]}, acento, line width=0.7pt] (0.6,-2.48) -- (0.6,-3.02);
  \node[draw=acento, rounded corners=3pt, anchor=north west, align=left, inner sep=5pt, line width=0.6pt] at (0,-3.05)
       {\textbf{363 latidos medios}\enspace{\footnotesize\color{gristexto}uno por paciente, $70 \times 12$ muestras}};
  \node[anchor=east] at (-0.2,-3.42) {Mediana por paciente};
\end{tikzpicture}}
\caption[Embudo de selección de latidos]{Embudo de selección de latidos con el criterio por defecto. La mayor pérdida (\SI{3.7}{\percent}) se debe a los latidos cortados por el inicio o el final del registro, no al control de calidad. Solo el \SI{0.9}{\percent} de los latidos se descarta por su morfología.}
\label{fig:embudo}
\end{figure}

\begin{table}[H]
\centering
\caption[Comparación de criterios de selección de instancias]{Comparación de criterios de selección de instancias. Cambio del latido medio: percentil 95, entre pacientes, de la máxima diferencia absoluta en V1--V2 respecto a no rechazar ningún latido. AUROC: regresión logística sobre las 72 características (media $\pm$ DT en 15 \emph{folds}).}
\label{tab:instancias}
\small
\begin{tabular}{@{}lS[table-format=3.0]S[table-format=1.1]S[table-format=2.0]S[table-format=1.3]S[table-format=1.3]@{\,$\pm$\,}S[table-format=1.3]@{}}
\toprule
\textbf{Criterio} & {\textbf{Descartados}} & {\textbf{\%}} & {\textbf{Pacientes}} & {\textbf{Cambio (mV)}} & \multicolumn{2}{c}{\textbf{AUROC}} \\
\midrule
Ninguno & 0 & 0.0 & 0 & {---} & 0.918 & 0.051 \\
$\rho < 0{,}5$ & 27 & 0.6 & 16 & 0.000 & 0.920 & 0.052 \\
\textbf{$\boldsymbol{\rho < 0{,}6}$ (elegido)} & 41 & 0.9 & 19 & 0.005 & 0.919 & 0.052 \\
$\rho < 0{,}7$ & 49 & 1.0 & 20 & 0.007 & 0.919 & 0.052 \\
$\rho < 0{,}8$ & 74 & 1.5 & 24 & 0.021 & 0.918 & 0.051 \\
$\rho < 0{,}9$ & 119 & 2.5 & 64 & 0.073 & 0.919 & 0.051 \\
RR desviado $>$\,\SI{20}{\percent} & 61 & 1.3 & 26 & 0.026 & 0.919 & 0.052 \\
$\rho < 0{,}6$ y RR & 87 & 1.8 & 35 & 0.047 & 0.920 & 0.052 \\
\bottomrule
\end{tabular}
\end{table}

\figura[\linewidth]{seleccion_instancias}{(a) Porcentaje de latidos descartados por cada criterio (barras) y AUROC del clasificador de referencia (puntos, $\pm$ error típico). Ningún criterio modifica el rendimiento. (b) Correlación del peor latido de cada paciente con su plantilla: el \SI{81.5}{\percent} de los pacientes no tiene ningún latido por debajo de \num{0.9}, en proporción similar en ambas clases.}{fig:instancias}

Vemos que el rechazo afecta por igual a ambas clases: descartamos el \SI{0.9}{\percent} de los latidos de los pacientes normales y un \SI{0.8}{\percent} de los de los pacientes que tienen Brugada. Por lo tanto, esta selección no introduce ningún tipo de sesgo que sea dependiente de la clase, las trata por igual. De hecho, como se ve en la \cref{fig:instancias}, los latidos de los pacientes que tienen Brugada son algo \emph{más} homogéneos que los de los pacientes sanos, teniendo una correlación mediana del peor latido de \num{0.951} en los Brugada, frente a \num{0.934} en los pacientes sanos (Mann-Whitney, $p = \num{0.001}$), por lo que con esto descartamos también que la clase positiva tenga una peor calidad de señal. Tenemos también que los criterios más estrictos descartan más latidos, pero esto apenas cambia el latido medio (como mucho \SI{0.07}{\milli\volt} en el percentil 95), porque la mediana ya ignoraba estos latidos.

Vemos que la selección de instancias no mejora el rendimiento en este conjunto de datos: los ocho criterios nos dan un AUROC entre \num{0.918} y \num{0.920}. No es un resultado negativo, sino que nos confirma dos cosas: que las señales, como dijimos antes, son limpias, y que la mediana ignora los latidos anómalos. Aun con esto, el criterio se va a mantener, porque no deja de ser una medida de protección frente a registros ruidosos, y deja constancia de lo que se descarta y de por qué se descarta.

\subsection{El latido medio por clase}

En la \cref{fig:latidoclase} se compara el latido medio de los grupos. Vemos que en los pacientes con Brugada, V1 y V2 muestran una onda r' más prominente, un segmento ST el cual desciende tras llegar al punto J en lugar de quedarse plano o incluso ascender, y una onda T de menor amplitud. Tenemos que la onda T es negativa en V2 en el \SI{43}{\percent} de los pacientes con Brugada, y solo en el \SI{6}{\percent} de los pacientes sanos. Este es el patrón tipo 1, que reconocemos en la \cref{fig:patrones}, aunque más atenuado, ya que se ha promediado con otros pacientes que no lo expresan en un grado tan alto. Vemos también que las bandas intercuartílicas se solapan ampliamente, que ninguna muestra aislada del latido separa las clases, y por eso necesitamos combinar varias medidas.

\figura[\linewidth]{latido_medio_clases}{Latido medio de V1 (a) y V2 (b) por clase: mediana entre pacientes (línea) y rango intercuartílico (banda). Las bandas grises marcan las regiones del QRS, del segmento ST (más oscura) y de la onda T.}{fig:latidoclase}

% ---------------------------------------------------------------------
\chapter{Extracción y selección de características}\label{cap:caracteristicas}

En este capítulo convertimos cada latido medio en un vector de 72 variables que aporten un significado clínico, las analizamos y se comparan cuatro métodos para seleccionar solo un subconjunto reducido de ellas.

\section{Extracción de características}

\subsection{Límites del QRS}

Casi todas las medidas dependen de dónde empieza y dónde acaba el QRS, porque es este el que define dónde se encuentra el punto J. Los delineadores habituales están pensados para trabajar entre 250 y \SI{1000}{\hertz} \cite{martinez2004}, por lo que no nos sirven para trabajar a los \SI{100}{\hertz} que tenemos nosotros. Por eso usamos un método, el cual está basado en la energía de la derivada (la pendiente de la señal) combinada de cuatro derivaciones complementarias: V1, V2, II y V5.

El QRS es la región de alta pendiente cerca del pico R. Su inicio y su final están situados donde esta energía cae por debajo del \SI{8}{\percent} de su máximo durante al menos 3 muestras consecutivas. Para evitar absurdos, el resultado se limita a rangos que sean verdaderamente posibles: el inicio entre \num{20} y \SI{120}{\milli\second} antes del pico R, y el final entre \num{20} y \SI{180}{\milli\second} después.

\subsection{Variables}

A partir de estos límites podemos definir 72 variables, organizadas en cuatro grupos (el catálogo completo está en el \cref{ap:catalogo}):

\begin{itemize}
  \item \textbf{Globales}: 3 en total, el inicio, el fin y la duración del QRS.
  \item \textbf{Por derivación, en V1, V2, V3 y II}: 16 por derivación, con un total de 64 variables.
  \begin{itemize}
    \item \emph{Segmento ST}: amplitud en el punto J y a \SIlist{40;80}{\milli\second}; pendiente, área y máximo del ST.
    \item \emph{Onda T}: amplitud, posición, mínimo e indicador de negatividad.
    \item \emph{Complejo QRS}: amplitudes de R, S y r'; área del QRS y pendiente final del QRS.
    \item \emph{Relativas}: cociente entre la elevación del ST y la amplitud del QRS.
  \end{itemize}
  \item \textbf{Agregados V1--V2}: 3 en total, máxima elevación del ST, mínimo de la onda T y máximo cociente ST/R.
  \item \textbf{Ritmo}: 2 en total, frecuencia cardíaca y coeficiente de variación del intervalo RR.
\end{itemize}

Como vemos, la pendiente del ST resulta ser la variable más importante de todas, y se define entre el punto J y \SI{80}{\milli\second} después. Que esta sea negativa indica un ST descendente, que es justamente el patrón tipo 1 (\cref{fig:patrones}). También incluimos la V3 porque el patrón puede extenderse a ella, y la II como derivación de control, alejada del ventrículo derecho.

\section{Análisis univariante}

En la \cref{tab:univariante} se nos muestran las variables más discriminantes cuando se evalúan de una en una, y en la \cref{fig:caracteristicas} se representa su distribución por clase. El AUC se presenta orientado, de modo que valga lo mismo para una variable la cual aumenta con la enfermedad que para una que disminuye.

\begin{table}[htbp]
\centering
\caption[Variables más discriminantes]{Variables más discriminantes en el análisis univariante (se omiten las casi equivalentes a otra de la tabla, como el final del QRS). Medianas por clase, AUC orientado y valor $p$ del test de Mann-Whitney.}
\label{tab:univariante}
\small
\begin{tabular}{@{}llS[table-format=-3.4]S[table-format=-3.4]S[table-format=1.3]l@{}}
\toprule
\textbf{Variable} & \textbf{Unidad} & {\textbf{Normal}} & {\textbf{Brugada}} & {\textbf{AUC}} & $\boldsymbol{p}$ \\
\midrule
Pendiente del ST, V1 & \si{\milli\volt\per\milli\second} & -0.0005 & -0.0012 & 0.845 & $2 \cdot 10^{-20}$ \\
Pendiente del ST, V2 & \si{\milli\volt\per\milli\second} & 0.0003 & -0.0008 & 0.809 & $1 \cdot 10^{-16}$ \\
Duración del QRS & \si{\milli\second} & 100 & 120 & 0.772 & $1 \cdot 10^{-13}$ \\
Amplitud de la onda T, V2 & \si{\milli\volt} & 0.339 & 0.183 & 0.720 & $3 \cdot 10^{-9}$ \\
Amplitud de r', V2 & \si{\milli\volt} & 0.198 & 0.308 & 0.719 & $4 \cdot 10^{-9}$ \\
Pendiente final del QRS, II & \si{\milli\volt\per\milli\second} & -0.0149 & -0.0088 & 0.707 & $3 \cdot 10^{-8}$ \\
Amplitud en el punto J, V2 & \si{\milli\volt} & 0.156 & 0.261 & 0.689 & $4 \cdot 10^{-7}$ \\
\midrule
ST a J+\SI{40}{\milli\second}, V2 & \si{\milli\volt} & 0.144 & 0.171 & 0.575 & \num{0.046} \\
ST a J+\SI{40}{\milli\second}, V1 & \si{\milli\volt} & 0.060 & 0.067 & 0.530 & \num{0.42} \\
\bottomrule
\end{tabular}
\end{table}

\begin{figure}[htbp]
\centering
\includegraphics[width=\linewidth]{caracteristicas_clave}
\caption[Distribución de las variables clave por clase]{Distribución por clase de cinco de las variables más discriminantes (a--e) y de la elevación del ST a J+\SI{40}{\milli\second} en V1 (f), que es la medida más próxima al criterio clínico de \SI{0.2}{\milli\volt} y apenas discrimina. En (c) se aprecian los escalones de \SI{10}{\milli\second} de la duración del QRS. Cajas: mediana y rango intercuartílico; puntos: pacientes individuales.}
\label{fig:caracteristicas}
\end{figure}

Lo que distingue a los pacientes que tienen Brugada no es cuánto se eleva el ST, sino la forma en la que este desciende. La pendiente del ST en V1 llega a alcanzar un AUC de \num{0.845}, mientras que la elevación del ST en J+\SI{40}{\milli\second} en la misma derivación se queda en \num{0.530} y no es significativa. Esto encaja con que la mayoría de los pacientes no lleguen a mostrar el patrón tipo 1 completo en el ECG basal, pero sí que muestren una tendencia descendente del ST, y que tengan los criterios de forma que están propuestos para el patrón tipo 2 \cite{chevallier2011,serra2014}.

\section{Redundancia}

Muchas variables miden aspectos que son parecidos del mismo latido, como el ST a \num{40} y a \SI{80}{\milli\second} dentro de una misma derivación, o tener la misma medida tanto para V1 como para V2. En la \cref{fig:correlacion} se muestra la matriz de correlación de las 72 variables, donde vemos que 30 pares tienen una correlación absoluta mayor de \num{0.9}, y 28 variables forman parte de al menos uno de esos pares. El hecho de que exista esta redundancia ya justifica tener que estudiar la selección de características. Y también nos enseña un riesgo, y es que los métodos que evalúan las variables de una en una pueden llegar a elegir varias que sean casi equivalentes.

\begin{figure}[htbp]
\centering
\includegraphics[width=0.94\linewidth]{correlacion_caracteristicas}
\caption[Matriz de correlación de las 72 características]{Matriz de correlación de Pearson entre las 72 características, ordenadas por agrupamiento jerárquico. Los bloques intensos en la diagonal agrupan medidas redundantes de una misma onda en una misma derivación.}
\label{fig:correlacion}
\end{figure}

\section{Métodos de selección comparados}

Podemos agrupar los métodos de selección de características dentro de tres familias \cite{guyon2003}: los \textbf{filtros}, que puntúan cada variable con un criterio estadístico sin la necesidad de usar ningún modelo; los \textbf{\eng{wrappers}}, que lo que hacen es probar subconjuntos de variables entrenando el modelo una y otra vez; y por último los \textbf{embebidos}, en los que la selección ocurre durante el propio entrenamiento. Se comparan dos filtros y dos embebidos, implementados mediante scikit-learn \cite{pedregosa2011}. Los \eng{wrappers} los dejo fuera, ya que al tener tan pocos registros de pacientes corremos el riesgo de que se sobreajusten.

\begin{itemize}
  \item \textbf{ANOVA F}: compara cuánto se separan las medias de las dos clases con cuánto varían los pacientes dentro de cada clase. Cuanto mayor es ese cociente, más útil es la variable.
  \item \textbf{Información mutua} \cite{kraskov2004}: mide cuánto se reduce la incertidumbre sobre la clase al conocer la variable. Esta es capaz de reconocer también relaciones no lineales.
  \item \textbf{Regularización L1} \cite{tibshirani1996}: es una regresión logística con penalización, que va arrastrando a cero los coeficientes de las variables que no son útiles.
  \item \textbf{Random Forest} \cite{breiman2001}: tenemos 500 árboles de decisión, y la importancia de cada variable se decide en función de cuánto mejoran, en promedio, los cortes que se hacen con ella.
  \item \textbf{Orden aleatorio}: usamos este como referencia, para ver cuánto aporta el seleccionar las variables frente a elegirlas al azar.
\end{itemize}

Vemos que todos estos métodos lo que producen es una ordenación de las 72 variables, de la más útil a la menos útil, lo que permite evaluarlos para cualquier número $k$ de variables.

Hay que tener cuidado, ya que si la ordenación se calculase con todos los pacientes y después evaluáramos con validación cruzada, las variables se habrían elegido viendo también el conjunto de pacientes de test, y se estaría trampeando el rendimiento \cite{ambroise2002}. Por eso, en cada uno de los 15 \eng{folds} (\cref{sec:protocolo}), cada método calcula su ordenación solo con el conjunto de entrenamiento; después se entrena una regresión logística con las $k$ primeras variables, para distintos valores de $k$ entre 1 y 72, y se mide el AUROC en el conjunto de test.

\section{Resultados}

Tenemos en la \cref{fig:fs-curvas} el rendimiento de cada método según el número de variables que usa, y en la \cref{tab:fs} se resumen los valores más relevantes.

\begin{figure}[htbp]
\centering
\includegraphics[width=0.75\linewidth]{seleccion_caracteristicas_curvas}
\caption[Rendimiento según el número de variables]{AUROC medio en 15 \eng{folds} según el número de variables seleccionadas (escala logarítmica). La línea de puntos marca el rendimiento con las 72 variables. Azul: filtros; verde: embebidos; gris: orden aleatorio.}
\label{fig:fs-curvas}
\end{figure}

\begin{table}[htbp]
\centering
\caption[Comparación de métodos de selección de características]{Comparación de métodos de selección de características: AUROC medio (15 \eng{folds}) con 3, 5 y 10 variables y tiempo medio por \eng{fold} (aproximado: depende del ordenador). Con las 72 variables todos coinciden: \num{0.919} $\pm$ \num{0.052}.}
\label{tab:fs}
\small
\begin{tabular}{@{}llS[table-format=1.3]S[table-format=1.3]S[table-format=1.3]S[table-format=1.2]@{}}
\toprule
 & & \multicolumn{3}{c}{\textbf{AUROC}} & \\
\cmidrule(lr){3-5}
\textbf{Método} & \textbf{Familia} & {$k=3$} & {$k=5$} & {$k=10$} & {\textbf{Tiempo (s)}} \\
\midrule
Regularización L1 & Embebido & 0.896 & 0.907 & 0.909 & {$\approx$\,0,5} \\
ANOVA F & Filtro & 0.888 & 0.904 & 0.901 & {$<$\,0,01} \\
Información mutua & Filtro & 0.893 & 0.900 & 0.902 & {$\approx$\,0,2} \\
Random Forest & Embebido & 0.893 & 0.898 & 0.905 & {$\approx$\,1,3} \\
\midrule
Orden aleatorio & Referencia & 0.688 & 0.727 & 0.780 & {---} \\
\bottomrule
\end{tabular}
\end{table}

En la \cref{fig:fs-frecuencia} se muestra con qué frecuencia cada método sitúa a cada variable dentro de las 10 primeras. Vemos que los 4 métodos coinciden en 5 variables casi siempre: la pendiente del ST en V1, la duración del QRS, la pendiente del ST en V2, la amplitud de la onda T en V2 y la amplitud en el punto J en V2. Todas estas describen el patrón característico del Brugada en las derivaciones precordiales derechas: el ST descendente, la T aplanada o negativa, la elevación del punto J y el retraso de conducción.

\begin{figure}[htbp]
\centering
\includegraphics[width=\linewidth]{seleccion_caracteristicas_frecuencia}
\caption[Frecuencia de selección por método]{Porcentaje de \eng{folds} en que cada método sitúa cada variable entre las 10 primeras, para las 12 variables más elegidas en promedio.}
\label{fig:fs-frecuencia}
\end{figure}

\section{Discusión}

Podemos sacar en claro tres cosas:

\begin{itemize}
  \item \textbf{La selección importa cuando vayamos a usar pocas variables}, ya que con 5 variables cualquier método supera al orden aleatorio en unos \num{0.17} o \num{0.18} de AUROC. Mientras que a partir de las 30 variables todos empiezan a converger, y ninguno mejora el usar las 72 variables.
  \item \textbf{Con 5 variables nos bastan}, ya que con la regularización L1 alcanzamos un AUROC de \num{0.907}, que es casi el \SI{99}{\percent} del rendimiento que alcanzamos con las 72 variables. Además, un modelo que tenga solo 5 medidas clínicas es mucho más fácil de explicar y de validar que uno que tenga 72.
  \item \textbf{Los cuatro métodos rinden prácticamente igual}: vemos que la diferencia que hay entre ellos es mucho más pequeña que la variación que hay entre los \eng{folds}, por lo que no se puede asegurar que uno sea mejor o peor. Lo que los diferencia de verdad es el coste, ya que el ANOVA es prácticamente instantáneo, mientras que el Random Forest es mucho más lento. En este problema funcionan tan bien los filtros sencillos porque las variables discriminantes ya lo son por separado.
\end{itemize}

% ---------------------------------------------------------------------
\chapter{Visualización mediante reducción de dimensionalidad}\label{cap:visualizacion}

La visualización nos permite comprobar si las clases efectivamente forman grupos separados, si hay pacientes atípicos y qué estructura tienen los datos. Al tener 72 variables no se pueden representar directamente, por lo que hay que proyectar los datos a dos dimensiones. Dentro de este capítulo voy a comparar tres técnicas.

\section{Métodos}

\paragraph{PCA.} Es el análisis de componentes principales \cite{jolliffe2002}, que busca las direcciones en las que los pacientes varían más y proyecta los datos sobre ellas. La primera componente es la dirección donde la varianza es máxima, la siguiente sería la de mayor varianza perpendicular a esta, y así sucesivamente. Es lineal y siempre nos devuelve el mismo resultado, y sus cargas nos permiten interpretar qué representa cada eje.

\paragraph{t-SNE.} Es un método no lineal \cite{vandermaaten2008} que intenta que los pacientes que son vecinos en el espacio original sigan siendo vecinos en el mapa. Su parámetro principal es la perplejidad, que regula a cuántos vecinos atiende cada punto. En mi caso usaré una perplejidad de 30.

\paragraph{UMAP.} Es un método que busca lo mismo que t-SNE \cite{mcinnes2018}: construye un grafo con los vecinos más cercanos a cada paciente y busca un mapa que sea capaz de conservarlos. Suele mantener mejor la estructura global y es más rápido que t-SNE. En mi caso usaré 15 vecinos.

En t-SNE y UMAP las distancias entre grupos y la forma global del mapa no tienen ningún significado y cambian según los parámetros \cite{wattenberg2016}; lo único fiable es quién está cerca de quién. Para no sacar conclusiones solo de una imagen, se comprueba además cuánta información sobre la clase conserva cada mapa: se entrena un $k$-NN usando solamente las dos coordenadas del mapa y se mide su AUROC.

Las tres técnicas se aplicaron sin usar etiquetas a dos representaciones estandarizadas: las 72 características y el latido medio de V1 y V2 muestra a muestra.

\section{Resultados}

La \cref{fig:dimred} muestra las seis proyecciones.

\begin{figure}[htbp]
\centering
\includegraphics[width=\linewidth]{reduccion_dimensionalidad}
\caption[Proyecciones 2D con PCA, t-SNE y UMAP]{Proyecciones 2D de los 363 pacientes con PCA, t-SNE y UMAP, a partir de las 72 características (fila superior) y del latido medio de V1--V2 (fila inferior). Azul: normal; rojo: Brugada.}
\label{fig:dimred}
\end{figure}

\subsection{PCA: la varianza no es la clase}

Las dos primeras componentes son capaces de explicar el \SI{22.5}{\percent} y el \SI{14.0}{\percent} de la varianza, y hacen falta hasta 20 componentes para llegar al \SI{90}{\percent} (\cref{fig:pca}). Las cargas nos muestran por qué la proyección no es capaz de separar las clases: la PC1 la forman las medidas de nivel del ST en V2, mientras que la PC2 la forman las mismas medidas pero en la derivación II. Estas son las direcciones en las que más varían los pacientes, pero no son las que distinguen entre tener Brugada o no. Como vimos en el \cref{cap:caracteristicas}, las que lo distinguen son la pendiente del ST y la duración del QRS. Un $k$-NN entrenado solo con las dos primeras componentes saca un AUROC de \num{0.53}, que es prácticamente lo mismo que tirar una moneda al aire.

\begin{figure}[htbp]
\centering
\includegraphics[width=\linewidth]{pca_varianza_cargas}
\caption[Varianza explicada y cargas de la PCA]{(a) Varianza explicada por cada componente (barras) y acumulada (línea). (b, c) Las ocho variables con mayor carga en PC1 y PC2, dominadas por el nivel del segmento ST en V2 y en II, respectivamente.}
\label{fig:pca}
\end{figure}

\subsection{t-SNE y UMAP: estructura local}

t-SNE y UMAP forman agrupaciones de pacientes con Brugada en los bordes del mapa. Sin embargo, esas agrupaciones están mezcladas con los pacientes normales, y el $k$-NN entrenado sobre el mapa (AUROC entre \num{0.71} y \num{0.76}) queda por debajo del que se obtiene usando todas las variables (\num{0.87} con las 72 características). Al cambiar la perplejidad o el número de vecinos, la forma de los mapas cambia, pero la mezcla de las clases se sigue manteniendo.

\section{Observaciones}

Como vimos, ninguna de las proyecciones a dos dimensiones es capaz de separar las clases. La PCA muestra que las direcciones de máxima varianza no se corresponden con las discriminantes, y t-SNE y UMAP solo nos dan agrupaciones locales parciales. Al estar la información repartida entre varias variables, nos hace falta un modelo supervisado que sea capaz de ponderarlas.

% ---------------------------------------------------------------------
\chapter{Modelo de clasificación de referencia}\label{cap:modelo}

En este último paso, lo que hacemos es comprobar si los datos que hemos preparado tienen un valor real, probándolo con un modelo sencillo. El objetivo ahora no es obtener el mejor clasificador posible, eso queda para el alcance del TFG, sino fijar una referencia de cuánto se puede llegar a discriminar con 72 características y un modelo lineal, y cuánto supera eso a aplicar directamente los criterios clínicos.

\section{Modelos}

\paragraph{Análisis discriminante lineal (LDA).} Este modelo supone que en cada clase las variables siguen una distribución normal, con la misma forma en las dos clases, y traza una frontera lineal que trate de separar de la mejor forma posible sus medias, teniendo en cuenta cómo varían y cómo se correlacionan las variables. Con las 72 variables muy correlacionadas (\cref{fig:correlacion}) y unos 290 pacientes de entrenamiento, estimar estas correlaciones es inestable. Por eso se usa la versión con contracción de Ledoit y Wolf \cite{ledoit2004}, que acerca la estimación a una más simple y calcula automáticamente cuánto hacerlo, sin tener que ajustar ningún hiperparámetro.

\paragraph{Regresión logística.} Combina linealmente las variables y transforma el resultado en una probabilidad de tener Brugada. Se usa con regularización L2, eligiendo su intensidad por validación cruzada interna, y con pesos de clase balanceados.

\paragraph{Reglas clínicas.} Se formalizan los tres criterios que caracterizan el patrón de Brugada con sus umbrales de consenso en V1 o V2: QRS $\geq$ \SI{120}{\milli\second}, ST $\geq$ \SI{0.2}{\milli\volt} en el punto J o a J+\SI{40}{\milli\second}, y onda T $\leq$ \SI{-0.1}{\milli\volt}. La regla «3 de 3» exige los tres criterios y se aproxima al patrón tipo 1. La regla «2 de 3» exige dos cualesquiera. Además, se define una puntuación continua, la máxima elevación del ST más la profundidad de la onda T negativa, con un umbral aprendido en el entrenamiento.

Todos los modelos se evalúan con el protocolo de la \cref{sec:protocolo}.

\section{Resultados}

\begin{table}[htbp]
\centering
\caption[Rendimiento del modelo de referencia]{Rendimiento del modelo de referencia frente a las reglas clínicas. AUROC: media $\pm$ DT en 15 \eng{folds} e IC 95\,\% \eng{bootstrap}. Las métricas con umbral usan el umbral de Youden elegido en el entrenamiento. Sens.: sensibilidad; Esp.: especificidad; VPP: valor predictivo positivo; VPN: valor predictivo negativo; BA: exactitud balanceada; MCC: coeficiente de correlación de Matthews.}
\label{tab:modelo}
\small
\setlength{\tabcolsep}{4.2pt}
\begin{tabular}{@{}lccS[table-format=1.2]S[table-format=1.2]S[table-format=1.2]S[table-format=1.2]S[table-format=1.2]S[table-format=1.2]@{}}
\toprule
\textbf{Modelo} & \textbf{AUROC} & \textbf{IC 95\,\%} & {\textbf{Sens.}} & {\textbf{Esp.}} & {\textbf{VPP}} & {\textbf{VPN}} & {\textbf{BA}} & {\textbf{MCC}} \\
\midrule
\textbf{LDA con contracción} & \textbf{\num{0.928} $\pm$ \num{0.047}} & \numrange{0.88}{0.96} & 0.87 & 0.88 & 0.67 & 0.96 & 0.88 & 0.69 \\
Regresión logística & \num{0.919} $\pm$ \num{0.052} & \numrange{0.86}{0.95} & 0.86 & 0.90 & 0.69 & 0.96 & 0.88 & 0.70 \\
\midrule
Puntuación ST + $|$T$|$ & \num{0.754} $\pm$ \num{0.053} & \numrange{0.69}{0.81} & 0.62 & 0.75 & 0.40 & 0.88 & 0.69 & 0.33 \\
Regla «2 de 3» & {---} & {---} & 0.78 & 0.71 & 0.42 & 0.92 & 0.75 & 0.41 \\
Regla «3 de 3» (tipo 1) & {---} & {---} & 0.22 & 0.99 & 0.80 & 0.83 & 0.60 & 0.36 \\
\bottomrule
\end{tabular}
\end{table}

En la \cref{tab:modelo} y en la \cref{fig:modelo} se recogen los resultados. El LDA llega a alcanzar un AUROC de \num{0.928}, con sensibilidad \num{0.87} y especificidad \num{0.88}. Supera a la puntuación clínica, que saca \num{0.754}, y la diferencia es mucho mayor que la variación que existe entre los diferentes \eng{folds}. Entre LDA y regresión logística tampoco vemos que difieran de una forma apreciable.

\begin{figure}[htbp]
\centering
\includegraphics[width=\linewidth]{modelo_referencia}
\caption[Curvas ROC y matriz de confusión]{(a) Curvas ROC de las predicciones fuera de muestra (primera repetición de la validación cruzada) para los modelos lineales y la puntuación clínica, con los puntos de operación de las reglas «2 de 3» y «3 de 3». (b) Matriz de confusión del LDA con su umbral de Youden, media de las tres repeticiones.}
\label{fig:modelo}
\end{figure}

Vemos cómo las reglas clínicas muestran el compromiso esperable. La regla «3 de 3» es casi perfecta en cuanto a especificidad (\num{0.99}), pero solo es capaz de detectar al \SI{22}{\percent} de los pacientes. Encaja con que la mayoría de los positivos no muestren el patrón completo en su ECG basal. Si la relajamos a «2 de 3», sube la sensibilidad a \num{0.78}, pero en cambio la especificidad baja a \num{0.71}. El modelo lineal es capaz de mejorar ambas a la vez: con las mismas variables, ponderarlas a partir de los datos supera a combinarlas sin más con umbrales fijos.

Además, tenemos también que el valor predictivo negativo del LDA es de \num{0.96}, por lo que un resultado negativo del modelo nos permite descartar el síndrome con bastante seguridad, mientras que uno positivo indica que conviene hacerle un estudio adicional.

A partir de esto, ya sabemos que con 72 características y un modelo lineal somos capaces de alcanzar un AUROC de \num{0.928}, el cual está muy por encima de los criterios clínicos. Esta es la cifra que debemos tratar de superar usando modelos más complejos dentro del TFG.

% ---------------------------------------------------------------------
\chapter{Discusiones y conclusiones}\label{cap:conclusiones}

\section{Conclusiones}

En este tipo de problema, la preparación de los datos es la que decide el resultado: vemos cómo unas pocas variables clínicas que estén bien construidas, usando un modelo lineal, son capaces de alcanzar un AUROC de entre \num{0.91} y \num{0.93}.

Por objetivos, hemos concluido que:

\begin{enumerate}[label=\textbf{O\arabic*}, leftmargin=2.4em]
  \item \textbf{ETL y calidad.} El conjunto de datos es de buena calidad: no contiene ningún valor ausente, tanto el ruido como la deriva son pequeños y todos los registros nos permiten construir un latido medio fiable. Durante el control de calidad fuimos capaces de detectar un registro duplicado, que podría haber causado una fuga entre el entrenamiento y el test.
  \item \textbf{Preprocesado.} El filtro elegido apenas influye, mientras que el detector de latidos sí que lo hace. En una señal muestreada a \SI{100}{\hertz}, la precisión temporal de la segmentación importa más que cómo de limpia esté la señal.
  \item \textbf{Selección de instancias.} Solo se descarta un \SI{0.9}{\percent} de los latidos, sin que exista un sesgo entre clases, y ninguno de los ocho criterios evaluados cambia el resultado. La mediana ya actúa como un estimador robusto, y la selección se mantiene como salvaguarda y por trazabilidad.
  \item \textbf{Características.} Sacamos en claro que la pendiente del segmento ST y la duración del QRS son capaces de discriminar mucho más que la elevación absoluta del ST, lo que coincide con los criterios de forma de la literatura clínica, como el índice de Corrado \cite{zorzi2015}. Los cuatro métodos de selección comparados rinden prácticamente igual, y coinciden en que con apenas 5 variables se conserva casi el \SI{99}{\percent} del rendimiento, así que bastaría con usar uno sencillo y barato, como puede ser el ANOVA.
  \item \textbf{Visualización.} Al final, ninguna de las proyecciones logró separar las dos clases. La PCA muestra que las direcciones de máxima varianza no son las discriminantes, y t-SNE y UMAP solo nos dejan ver agrupaciones locales parciales.
  \item \textbf{Modelo de referencia.} El LDA con contracción alcanza un AUROC de \num{0.928} (IC 95\,\%: \numrange{0.88}{0.96}), un valor muy superior al obtenido con la puntuación basada en los criterios clínicos (\num{0.75}).
\end{enumerate}

\section{Limitaciones}

\begin{itemize}
  \item Al trabajar a \SI{100}{\hertz}, cada muestra equivale a \SI{10}{\milli\second}. Por eso la duración del QRS solo puede tomar valores escalonados, y el punto J se localiza con una incertidumbre de una muestra.
  \item Los datos que tenemos son solo de un centro y de un ECG por paciente, y no tenemos ningún tipo de validación externa de estos registros.
  \item El tamaño muestral es muy bajo también: contando solo con unos 15 positivos por \eng{fold} de test, las diferencias de AUROC que sean menores de \numrange{0.02}{0.03} no llegan a ser concluyentes.
\end{itemize}

\section{Continuación con el Trabajo de Fin de Grado}

El proceso y el conjunto de datos construidos aquí representan la base de la parte comparativa de mi TFG. Con el mismo protocolo de evaluación, en el TFG compararé siete familias de métodos: reglas clínicas, modelos clásicos, autoencoders, redes convolucionales y de atención sobre la señal, el modelo preentrenado ECG-FM \cite{mckeen2025}, redes sobre imágenes del latido y modelos multimodales de visión-lenguaje \cite{costa2026vlm}. También se analizarán la interpretabilidad de los modelos, la calibración de sus probabilidades y los pacientes que todos los métodos clasifican mal. El valor de referencia que fijamos en el \cref{cap:modelo} es la cifra que tratarán de superar todos estos nuevos y más complejos modelos.


\printbibliography[heading=bibintoc, title={Bibliografía}]

\appendix
\crefalias{chapter}{appendix}
% ---------------------------------------------------------------------
\clearpage
\section{Catálogo de características}\label{ap:catalogo}

Las 72 características se calculan sobre el latido medio de cada paciente (70 muestras por derivación, pico R en la muestra 25, nivel del segmento PR restado). Su implementación está en \var{src/brugada/features.py}.

\subsection*{Características globales y de ritmo (5)}

\begin{center}
\small
\begin{tabularx}{\linewidth}{@{}p{3.6cm}Lp{1.5cm}@{}}
\toprule
\textbf{Nombre} & \textbf{Definición} & \textbf{Unidad} \\
\midrule
\var{qrs\_onset\_ms} & Inicio del QRS respecto al pico R & ms \\
\var{qrs\_offset\_ms} & Final del QRS (punto J) respecto al pico R & ms \\
\var{qrs\_dur\_ms} & Duración del QRS & ms \\
\var{hr\_bpm} & Frecuencia cardíaca media del registro & lpm \\
\var{rr\_cv} & Coeficiente de variación de los intervalos RR & --- \\
\bottomrule
\end{tabularx}
\end{center}

\subsection*{Características por derivación (16 $\times$ 4 = 64)}

Se calculan en V1, V2, V3 y II, con el nombre de la derivación como prefijo (por ejemplo, \var{V1\_st\_slope}).

\begin{center}
\small
\begin{tabularx}{\linewidth}{@{}p{3.6cm}Lp{1.5cm}@{}}
\toprule
\textbf{Sufijo} & \textbf{Definición} & \textbf{Unidad} \\
\midrule
\var{st\_j} & Amplitud en el punto J (final del QRS) & mV \\
\var{st\_j40} & Amplitud a J+40 ms & mV \\
\var{st\_j80} & Amplitud a J+80 ms & mV \\
\var{st\_slope} & Pendiente del ST entre J y J+80 ms & mV/ms \\
\var{st\_area} & Área del ST entre J y J+80 ms & mV·ms \\
\var{st\_max} & Máximo del ST entre J y J+120 ms & mV \\
\var{t\_amp} & Amplitud con signo del extremo de la onda T (de J+100 ms al final) & mV \\
\var{t\_pos} & Posición del extremo de la onda T respecto al pico R & ms \\
\var{t\_neg} & Mínimo de la onda T & mV \\
\var{t\_negative} & Indicador de T negativa: mínimo $< -0{,}1$ mV y mayor en valor absoluto que la mitad del máximo & 0/1 \\
\var{r\_amp} & Máximo del QRS (onda R) & mV \\
\var{s\_amp} & Mínimo del QRS (onda S) & mV \\
\var{qrs\_area} & Área absoluta del QRS & mV·ms \\
\var{rprime\_amp} & Máximo del QRS posterior a la onda S (onda r') & mV \\
\var{qrs\_end\_slope} & Pendiente desde el extremo del QRS hasta el punto J (forma convexa) & mV/ms \\
\var{st\_r\_ratio} & Cociente entre la amplitud a J+40 ms y la máxima amplitud absoluta del QRS & --- \\
\bottomrule
\end{tabularx}
\end{center}

\subsection*{Agregados de V1 y V2 (3)}

\begin{center}
\small
\begin{tabularx}{\linewidth}{@{}p{3.6cm}Lp{1.5cm}@{}}
\toprule
\textbf{Nombre} & \textbf{Definición} & \textbf{Unidad} \\
\midrule
\var{v12\_st\_j40\_max} & Máximo de \var{st\_j40} entre V1 y V2 & mV \\
\var{v12\_t\_neg\_min} & Mínimo de \var{t\_neg} entre V1 y V2 & mV \\
\var{v12\_st\_r\_ratio\_max} & Máximo de \var{st\_r\_ratio} entre V1 y V2 & --- \\
\bottomrule
\end{tabularx}
\end{center}

Además se calculan dos indicadores de calidad (número de latidos conservados y correlación media con la plantilla), que se usan en el control de calidad pero no como variables predictivas.

% ---------------------------------------------------------------------
\clearpage
\section{Reproducibilidad}\label{ap:reproducibilidad}

Todo el trabajo se reproduce con el repositorio público \url{https://github.com/CesarCapelo/BrugadaInteligenciaDeNegocio}. El conjunto de datos se descarga automáticamente de PhysioNet. Las semillas aleatorias están fijadas (semilla 42, con un desplazamiento determinista por \emph{fold}), y la huella SHA-256 de los ficheros de señal (prefijo \var{d714832551a217e6}) garantiza que los resultados corresponden a la misma versión de los datos.

\subsection*{Entorno}

Python 3.12, NumPy 1.26.4, SciPy 1.13.1, scikit-learn 1.5.2, NeuroKit2 0.2.10, WFDB 4.3.1, umap-learn 0.5.7, pandas 2.2.3 y Matplotlib 3.9.4. La lista completa está en \var{requirements.txt}.

\subsection*{Comandos}

\begin{center}
\small
\begin{tabularx}{\linewidth}{@{}p{4.3cm}L@{}}
\toprule
\textbf{Análisis} & \textbf{Comando (desde \var{src/})} \\
\midrule
Descarga de los datos & \var{bash ../scripts/download\_data.sh} \\
ETL y control de calidad (\cref{cap:datos}) & \var{python -m brugada.experiments.run\_etl} \\
Comparación del preprocesado (\cref{cap:preprocesado}) & \var{python -m brugada.experiments.run\_preprocessing} \\
Selección de instancias (\cref{cap:instancias}) & \var{python -m brugada.experiments.run\_instance\_selection} \\
Características y su selección (\cref{cap:caracteristicas}) & \var{python -m brugada.experiments.run\_features} \\
Reducción de dimensionalidad (\cref{cap:visualizacion}) & \var{python -m brugada.experiments.run\_dimred} \\
Modelo de referencia (\cref{cap:modelo}) & \var{python -m brugada.experiments.run\_baseline} \\
Figuras de este documento & \var{python ../memoria/make\_figures.py} \\
\bottomrule
\end{tabularx}
\end{center}

Cada experimento guarda en \var{results/} sus métricas por \emph{fold}, el resumen y las predicciones fuera de muestra de cada paciente en cada repetición, de modo que cualquier tabla o figura puede regenerarse sin volver a entrenar.

% ---------------------------------------------------------------------
\chapter{Uso de herramientas de inteligencia artificial}\label{ap:ia}

En la elaboración de este trabajo se ha utilizado un asistente de inteligencia artificial generativa (Claude, de Anthropic) como apoyo en la programación del proceso de análisis, en la generación de figuras y en la redacción y revisión del texto, con el conocimiento y el permiso del profesor de la asignatura. El autor ha definido el planteamiento del trabajo, ha revisado y validado el código, los resultados y el texto, y asume la responsabilidad sobre su contenido. Todos los resultados se obtienen de código publicado y reproducible (\cref{ap:reproducibilidad}).

\end{document}
